\documentclass[10pt,a4paper,twocolumn]{article}
\usepackage[revtexheadings]{natureastro-custom}

\title{A Nature Astronomy-inspired manuscript in \LaTeX}
\author{First Author$^{1}$, Second Author$^{2}$ \& Third Author$^{1}$}
% \runninghead{First Author et al.}

\begin{abstract}
  This unofficial template is intended for writing manuscripts, notes and polished preprints in a visual style inspired by Nature Astronomy. It keeps the useful parts of the published layout---a full-width title and abstract, restrained sans-serif typography, compact two-column text, superscript numerical citations and concise captions---without production-only metadata such as received, accepted or publication dates, DOI information or volume details.
\end{abstract}

\affiliations{%
  $^{1}$School of Physics, University of Example, Melbourne, VIC, Australia.\\
$^{2}$Institute for Astronomy, Example University, Example City, Country.}
\correspondence{corresponding.author@example.edu}

\begin{document}
\maketitle

A manuscript can begin directly with the scientific narrative rather than with a heading labelled ``Introduction''. The opening paragraphs should establish the context, identify the unresolved problem and state the contribution succinctly. Citations appear as numerical superscripts in order of appearance.\cite{einstein1916,shakura1973}

This file is deliberately a clean authoring template rather than a simulation of the journal's production workflow. Its purpose is to make drafts pleasant to read while keeping the source conventional and easy to adapt later.

\section{Results}
\subsection{A compact two-column hierarchy}
The principal body is set in two columns with narrow gutters. Section headings are compact and bold. Displayed equations use conventional mathematical typography while the surrounding prose remains sans serif:
\begin{equation}
  \frac{\mathrm{d}\boldsymbol{x}}{\mathrm{d}t}
  = \boldsymbol{f}(\boldsymbol{x},t) + \boldsymbol{\eta}(t),
\end{equation}
where $\boldsymbol{\eta}$ denotes a stochastic forcing term.

\begin{figure}[t]
  \centering
  \fbox{\rule{0pt}{35mm}\rule{0.94\columnwidth}{0pt}}
  \caption{\textbf{Example single-column figure.} Replace the framed placeholder with \texttt{\textbackslash includegraphics}. Captions should explain symbols and uncertainties sufficiently for the figure to be interpreted without repeatedly returning to the main text.}
  \label{fig:example}
\end{figure}

Figure~\ref{fig:example} illustrates the default single-column treatment. For a figure spanning the whole page width, use the standard \texttt{figure*} environment.

\subsection{Tables}
\begin{table}[t]
  \caption{Example compact table}
  \centering
  \begin{tabular}{lcc}
    \toprule
    Quantity & Symbol & Value \\
    \midrule
    Spin frequency & $\nu$ & 100\,Hz \\
    Timescale & $\tau$ & 10\,d \\
    Fraction & $x$ & 0.10 \\
    \bottomrule
  \end{tabular}
\end{table}

\section{Discussion}
Use the layout as a writing environment for a manuscript, internal draft, research note or preprint. The source remains ordinary \LaTeX{}, so the manuscript can later be transferred to a journal-specific template without restructuring the content.

\methods
Methods can continue in the same two-column format. For a real manuscript, place sufficient implementation detail here to reproduce the analysis, together with statistical conventions, software versions and selection criteria.

\dataavailability
State where the data underlying the study can be obtained, if this section is useful for the document.

\codeavailability
State where analysis code is archived, if applicable.

\acknowledgements
Acknowledge funding, facilities, collaborations and other contributions as needed.

\authorcontributions
F.A. conceived the study. F.A. and S.A. performed the analysis. All authors discussed the results and contributed to the manuscript.

\competinginterests
The authors declare no competing interests.

\begin{thebibliography}{9}
  \bibitem{einstein1916}
  Einstein, A. The foundation of the general theory of relativity. \textit{Ann. Phys.} \textbf{49}, 769--822 (1916).

  \bibitem{shakura1973}
  Shakura, N. I. \& Sunyaev, R. A. Black holes in binary systems. Observational appearance. \textit{Astron. Astrophys.} \textbf{24}, 337--355 (1973).
\end{thebibliography}
% \bibliographystyle{aasjournalv7}

\makeendmatter
\end{document}
